% Generated from locale/tr/content/set-theory/ord-arithmetic/introduction.tex; do not hand-edit.

\olfileid[tr]{sth}{ord-arithmetic}{intro}

\olsection{Giriş}

\olref[ordinals][]{chap} içinde sıra sayıları teorisini geliştirdik.
\olref[spine][]{chap} içinde sıra sayılarını, hiyerarşinin geri kalanının
çevresinde kurulduğu bir omurga gibi düşünebileceğimizi gördük. Ne var ki
sıra sayılarının tek rolü bu değildir. Bir de sıra sayısı aritmetiğini
yapma görevi vardır.

Aslında \olref[ordinals][idea]{sec} içinde $\omega$, $\omega+1$ ve
$\omega+\omega$'dan söz ederken buna işaret etmiştik. O zaman gayriresmî
konuşmuştuk; şimdi ayrıntıları gerektiği gibi açıklamanın zamanı geldi. Yine
de bu bölümde fazla felsefe bulunmadığını belirtelim: yalnızca teknik
gelişmeler ve aynı şeyi iki farklı yoldan yapabileceğimize ilişkin (biraz)
ilginç bir gözlem vardır.

